% =====================================================================
% DTW: вывод рекуррентного соотношения (фундамент OLTW)
% =====================================================================
\begin{frame}{Dynamic Time Warping --- основа OLTW (определение)}

\textbf{Дано:} партитура $P = (p_1, p_2, \ldots, p_N)$ (последовательность
pitch-ов) и наблюдаемое исполнение $O = (o_1, o_2, \ldots, o_T)$.

\bigskip

\textbf{Ищем} \emph{warping path} $\pi = ((i_1, t_1), \ldots, (i_L, t_L))$
с тремя свойствами:

\begin{enumerate}
\item граничные условия: $(i_1, t_1) = (1, 1)$, $(i_L, t_L) = (N, T)$;
\item монотонность: $i_{\ell+1} \ge i_\ell$ и $t_{\ell+1} \ge t_\ell$;
\item шаг-1: $(i_{\ell+1} - i_\ell, t_{\ell+1} - t_\ell) \in \{(0,1),(1,0),(1,1)\}$.
\end{enumerate}

\bigskip

\textbf{Стоимость пути}:
\[
\mathrm{cost}(\pi) \;=\; \sum_{\ell=1}^{L} c(p_{i_\ell}, o_{t_\ell}),
\qquad
c(p, o) = |p - o| \text{ (полутоны)}.
\]

\bigskip

\textbf{Оптимальный путь} $\pi^* = \argmin_{\pi} \mathrm{cost}(\pi)$.

\end{frame}

\begin{frame}{DTW --- рекуррентное соотношение (вывод)}

Обозначим $D[i, t]$ --- минимальная стоимость пути из $(1,1)$ в $(i, t)$.

\bigskip

\textbf{Идея вывода:} последний шаг пути приходит в $(i, t)$ из одного
из трёх предков (свойство шаг-1):

\begin{itemize}
  \item диагональ $(i{-}1, t{-}1) \to (i, t)$;
  \item горизонталь $(i{-}1, t) \to (i, t)$;
  \item вертикаль $(i, t{-}1) \to (i, t)$.
\end{itemize}

\bigskip

\textbf{Принцип оптимальности Беллмана:} путь до $(i, t)$ оптимален
$\Longleftrightarrow$ путь до предыдущего шага оптимален. Отсюда:

\[
\boxed{\;
D[i, t] \;=\; c(p_i, o_t) \;+\;
\min\!\Bigl(D[i{-}1, t{-}1],\; D[i{-}1, t],\; D[i, t{-}1]\Bigr)\;}
\]

с базовым случаем $D[1,1] = c(p_1, o_1)$ и $D[\cdot, \cdot] = +\infty$
вне сетки.

\bigskip

\textbf{Сложность:} $O(NT)$ по памяти и времени.

\textbf{Online-проблема:} в realtime $T$ растёт, нельзя хранить весь
массив. \rlhi{Решение --- OLTW.}

\end{frame}
